\section{\label{sec:ope}算符乘积展开}

Wilson 处理非定域算符 OPE 的方法广为人知——例如可参见 Ref.~\cite{Collins:1984rdg}——
这里我们仅回顾讨论中必需的一些记号。

我们把非定域算符 ${\cal Q}(x)$ 的 OPE 写为
\begin{equation}
{\cal Q}(x)  \stackrel{x\rightarrow 0}{\sim} \; \sum_k
c_k(x,\mu) [{\cal O}^{(k)}(0)]_{\mathrm{R}}.
\end{equation}
威尔逊系数 $c_k(x,\mu)$ 是时空间隔 $x$ 与重正化尺度 $\mu$ 的复函数，承载与重正化局域算符
$[{\cal O}^{(k)}(0)]_{\mathrm{R}}$ 相联系的短程物理。我们用
$[\ldots]_{\mathrm{R}}$ 表示重正化算符，并出于记号简便省略其对重正化尺度 $\mu$ 的依赖。
局域算符的自由场质量维数决定相应威尔逊系数的带头时空依赖。

作为例子，考虑两个标量场的编时两点函数，时空间隔为 $x$：${\cal T} \{\phi(x)\phi(0)\}$。
在自由标量场论中，
OPE 是一个洛朗展开。相互作用在威尔逊系数中产生次领先的时空依赖，使其成为时空间隔、（重正化）
质量 $m_{\mathrm{R}}$ 与重正化尺度 $\mu$ 的函数：
\begin{align}\label{eq:opeint}
{\cal T} \{\phi(x)\phi(0)\} = {} & \frac{c_{\mathbb{I}}(\mu
x,m_{\mathrm{R}}x)}{4\pi^2x^2}\mathbb{I}+
c_{\phi^2}(\mu x,m_{\mathrm{R}}x)[\phi^2(0) ]_{\mathrm{R}} \nonumber\\
{} & \qquad
+{\cal O}(x).
\end{align}
这里我们已经把带头时空依赖从威尔逊系数中提出来。${\cal O}(x)$ 表示数量级为 $x$
的项，直至相互作用生成的对数修正。

我们建议把 OPE 中的那组局域算符替换为它们的局域涂抹对应物
\begin{equation}
{\cal Q}(x)  \stackrel{x\rightarrow 0}{\sim} \; \sum_k
d_k(\tau,x,\mu) {\cal S}^{(k)}(\tau,0).
\end{equation}
此时，威尔逊系数 $d_k(\tau,x,\mu)$ 与涂抹算符 ${\cal S}^{(k)}(\tau,0)$ 都是一个额外尺度——流时间
$\tau$——的函数。尽管如此，威尔逊系数的带头时空依赖仍由相应涂抹算符的正则质量维数决定，
因而与标准 OPE 的完全相同。

正如 OPE 仅在小时间空间隔下有效，sOPE 也要求小流时间。我们将看到，需要 $\tau \propto x^2$
（其中 $x$ 设定为小量），以确保威尔逊系数与外部态无关。

例如，回到编时两点函数，其 sOPE 为
\begin{align}\label{eq:sope2pt}
\!\!{\cal T}\{\phi(x)\phi(0)\} ={} &
\frac{d_{\mathbb{I}}(\mu x,\mu^2\tau,m_{\mathrm{R}}x)}{4\pi^2x^2}\mathbb{I}
\nonumber\\
{} & +
d_{\rho^2}(\mu x,\mu^2\tau,m_{\mathrm{R}}x)\rho^2(\tau,0)  +
{\cal O}(x,\tau),
\end{align}
其中我们把流时间 $\tau$ 处的涂抹场记作 $\rho(\tau,x)$。

在下面几节中，我们将观察到 sOPE 中涂抹威尔逊系数的四个特征。
\begin{enumerate}
\item 原 OPE 的对数时空依赖在 sOPE 中得以保持。
\item 在微扰论领头阶，流时间充当涂抹威尔逊系数的紫外调节子。
\item 超出领头阶，流时间无法再正则化威尔逊系数，因为流演化是经典的。不过，我们可以把重正化尺度
依赖吸收进原理论的重正化参数之中。
\item 对任何有意义的 OPE 而言，威尔逊系数都必须与外部态无关。对 sOPE，可以通过选取
$s_{\mathrm{rms}} < x$ 来保证这一点；也就是说，平均涂抹半径必须小于非定域算符的时空延展。
这保证 sOPE 仍是关于近似局域算符的展开。换言之，如果梯度流探测的长度尺度达到非定域算符的量级，
那么 sOPE 对原算符而言就是一个糟糕的展开。
\end{enumerate}

尽管我们把这个展开称为\emph{涂抹} OPE，但我们心目中的涂抹是通过梯度流实现的。
梯度流作为一种光滑化操作，把理论的自由度推向作用量的驻点。在 QCD 中，梯度流对应于连续的
stout 涂抹过程——一种构造紫外涨落受到阻尼的格点规范场的解析方法 \cite{Morningstar:2003gk}。
梯度流与其他涂抹方案的直接比较最早见于 Ref.~\cite{Bonati:2014tqa}。

许多梯度流研究对小流时间的场按零流时间的局域场做小流时间展开
\cite{Shindler:2013bia,Suzuki:2013gza,DelDebbio:2013zaa,
Makino:2014taa,Asakawa:2013laa,Luscher:2014kea}。我们可以把这样的展开视为流时间方向的 OPE，
从而把在非零流时间计算的重正化量与原理论中零流时间的相应量联系起来——否则后者难以计算。

本工作采取不同的思路。我们不把流动场按非零流时间的原场展开，而是把正流时间处的（矩阵元）场本身
作为基本研究对象。两种做法都反映了一个物理上有依据的预期：远小于系统物理尺度的涂抹尺度不应扭曲
所关心的物理。小流时间展开定量刻画了对这一预期的任何偏离，并且进一步把连续统中威尔逊系数的解析
计算与采用涂抹自由度的格点计算解耦。sOPE 则把涂抹尺度作为系统的内禀尺度纳入其中，这就要求确定新的
威尔逊系数。于是，小流时间展开与 sOPE 都由涂抹尺度作为紫外调节子的角色塑造，是对同一物理
（在格点上部分恢复转动对称性）既相互关联又彼此不同的概念性处理。

一般而言，涂抹是部分恢复转动对称性 \cite{Davoudi:2012ya} 从而改善格点计算连续极限的工具。
经由梯度流的涂抹具有若干优点。特别地，使用涂抹自由度在格点上非微扰确定的矩阵元无须进一步重正化
\cite{Luscher:2010iy}（费米子波函数重正化除外 \cite{Luscher:2013cpa}），
且只要在取连续极限时涂抹尺度以物理单位保持固定，便保持有限。

现在我们转向对 $\phi^4$ 标量场论中 sOPE 更完整的研究；这是审视 sOPE 时特别直截了当的一个理论。
由于流时间演化对标量场是线性的，我们可以精确求解流时间方程。
我们也无须考虑与规范固定相关的复杂之处
\cite{Luscher:2010iy,Luscher:2011bx}——以及费米子的额外重正化 \cite{Luscher:2013cpa}。
此外，对标量场而言，sOPE 可以理解为就是原 OPE 的一个重求和。虽然这对本文并非必需，但顺带指出也
颇有意味：已知在四维欧氏 $\phi^4$ 理论中，固定（但任意）微扰阶的 OPE 是收敛的
\cite{Hollands:2011gf}。该结果对任意时间空间隔均成立，只要外部态具有紧致支撑。

展望未来在 QCD 中的计算，我们预计含规范场相互作用的流方程所带来的技术问题会因增加给定微扰阶的图数
而略微加重微扰计算的负担，但不会改变我们的结论。在规范部分，不存在流动场的圈，因为在非零流时间处
重正化关联函数保持有限 \cite{Luscher:2011bx}。因此在微扰论领头阶，流时间正则化紫外发散；
超出领头阶则必须纳入适当的重正化手续。对费米子场，存在一个额外的费米子重正化参数，
已在文献 \cite{Luscher:2013cpa} 中计算，但它可以通过考虑重正化群不变量而移除
\cite{Luscher:2013cpa}。我们还注意到，sOPE 的所有微扰计算都可以在连续统中进行，
无须诉诸通常更为繁琐的格点微扰论 \cite{Monahan:2013dla,Capitani:2002mp}。
